\chapter{教学作为表演艺术}\label{s:performance}

在\emph{《机器的达尔文》}（\emph{Darwin Among the Machines}）中，
George Dyson\index{Dyson, George}写道：
「在生命与进化的牌桌上，有三位玩家：
人类、自然和机器。
我坚定地站在自然这一边。
但我怀疑，自然站在机器那一边{\ldots}」
如今教育这场牌桌上也类似地有三位玩家：
教科书和其他只读材料、
现场讲座，
以及自动化的在线课程。
你可以同时给你的学习者书面课程，
以及录制视频和自主进度练习的某种组合，
但如果你要当面授课，
你就必须提供某种不同于（并且希望能好于）这两者的东西。
因此本章的重点，是如何通过真正动手做来教编程。

\seclbl{现场编程}{s:performance-live}

\begin{quote}

  教学是戏剧，不是电影。\\
  —— Neal Davis\index{Davis, Neal}

\end{quote}

教编程最有效的方式是\gref{g:live-coding}{现场编程}~\cite{Rubi2013,Haar2017,Raj2018}。
老师不是呈现预先写好的材料，
而是在全班面前写代码，
学习者一边跟着敲、一边运行。
现场编程比幻灯片更好，原因有几个：

\begin{itemize}

\item
  它让\gref{g:active-teaching}{主动教学}成为可能，
  因为老师可以在当下顺着学习者的兴趣和问题走。
  一套幻灯片就像铁轨：
  坐上去也许很平稳，
  但你得提前决定去哪。
  现场编程更像开越野车：
  路可能更颠，
  但改变方向、去大家想去的地方要容易得多。

\item
  看一个程序被写出来，
  比看别人一页页翻幻灯片更有动力。

\item
  它促进了非预期知识传递：\index{非预期知识传递}
  通过观察我们\emph{怎么}做事，人们学到的东西比我们有意教的更多。

\item
  它让老师慢下来：
  如果要一边讲一边把程序敲进去，
  他们最多只能比学习者快一倍，
  而不像用幻灯片时那样能快上十倍。

\item
  它有助于减轻短期记忆的负担，
  因为它让老师更清楚自己往学习者身上抛了多少东西。

\item
  学习者能看到如何诊断和纠正错误。
  他们将来会花大量时间做这件事；
  除非你是完美的打字员，
  现场编程能确保他们看到该怎么做。

\item
  看老师犯错，能让学习者明白
  自己犯错也没关系。
  如果老师不因为犯错、谈起错误而难堪，
  学习者做这些事时也会更自在。

\end{itemize}

现场编程的另一个好处，是它展示了程序该按什么顺序写。
在研究人们如何解决帕森斯问题时，\index{帕森斯问题}
\cite{Ihan2011}发现，有经验的程序员常常先把方法签名拖到开头，
然后加上大部分控制流（也就是循环和条件），
最后才补上变量初始化、边界情况处理之类的细节。
这种乱序书写对新手很陌生，
因为他们是按代码在页面上出现的顺序读和写的；
看到这种写法，能帮他们学会把问题分解成可以逐个攻克的子目标。
现场编程还给了老师一个机会，去强调小步前进、频繁反馈的重要性~\cite{Blik2014}，
以及在改代码时先选好一个方案、
而不是多少有点随机地乱改、寄希望于碰巧变好~\cite{Spoh1985}。

一边在人前写代码一边说话，需要一点练习才能适应，
但大多数人反映，它很快就变得和对着幻灯片讲话一样不费力。
下面几节给出一些让现场编程变得更好的建议。

\subsection*{拥抱你的错误}

\begin{quote}

  笔误就是教学法。\\
  —— Emily Jane McTavish\index{McTavish, Emily Jane}

\end{quote}

现场编程最重要的规则，是拥抱你的错误。\index{错误（拥抱错误的重要性）}
不管你准备得多充分，
总还是会犯一些；
犯的时候，
就和你的听众一起把它们想清楚。
虽然很难找到确切数据，
但专业程序员要把 25\% 到 60\% 的时间花在调试上；
新手花得更多（\secref{s:pck-debug}），
然而大多数教科书和教程几乎不花时间讲如何诊断和纠正问题。
如果你在弄清自己哪里打错、或哪一步走偏时把想法说出来，
并解释自己是怎么纠正的，
你就给了学习者一套工具箱，让他们在自己犯错时也能用。

\begin{aside}{刻意的「手滑」}
  一门课讲过几遍之后，
  你除了基本的打字错误（这些仍然有启发性）之外，就不太会犯别的错了。
  你可以试着回忆过去的错误并故意再犯，
  但这往往显得很别扭。
  另一个办法是\gref{g:twitch-coding}{直播编程}：
  让学习者一个接一个地告诉你下一步该敲什么。
  这几乎保证会让你陷入某种麻烦。
\end{aside}

\subsection*{让大家做预测}

在一边现场编程时保持学习者投入的一个办法，
是让他们预测屏幕上的代码接下来会做什么。
你可以先记下他们给出的头几个猜测，
让全班投票看哪个最可能，
然后再把代码跑起来。
这是一种轻量的同伴教学，\index{同伴教学}
我们会在\secref{s:classroom-peer}里讨论；
它既能让他们的注意力保持在任务上，
又能让他们练习如何推理代码的行为。

\subsection*{慢一点}

每次你敲一条命令、
往程序里加一行代码、
或从菜单里选一项，
都把你正在做的事大声说出来，
然后指着屏幕上你做完的东西及其输出，
再从头过一遍。
这让学习者跟得上，
也能核对刚才做的对不对。
当有些学习者有视觉或听觉障碍、
或对教学语言不够流利时，这一点尤其重要。

不管你做什么，
\emph{都不要}复制粘贴代码：
这么做几乎保证你会把学习者远远甩在后面。
如果你用了 Tab 补全，
就把它说出来，好让学习者明白你在做什么：
「我们用 turtle 点『r』『i』再按 Tab，得到『right』。」

如果你敲的命令或代码的输出让刚敲的内容从视野里消失了，
就往上滚动，让学习者能重新看到它。
如果不方便这么做，
就把同一条命令再执行一遍，
或者把最后几条命令复制粘贴到工作坊的共享笔记里。

\subsection*{让人看得见、听得清}

你一坐下，
就更可能盯着自己的屏幕而不是看听众，
还可能被后排的学习者挡住。
如果你体力上能站上一两个小时，
那么讲课时就应该站着。
为此提前安排，确保你有升高的桌子、
站立式办公桌，
或讲台
来放笔记本电脑，
这样就不必弯着腰打字。

不管你站着还是坐着，
都要尽量多走动：
走到屏幕前指一指某个东西、
在白板上画点什么，
或者干脆离开电脑一会儿，直接对听众讲话。
这么做会把学习者的注意力从他们的屏幕上引开，
也给了他们一个自然的提问时机。

如果你要讲超过两个小时，
即便在小房间里也值得用麦克风。
你的嗓子会累，就像身体其他任何部位一样；
用麦克风和穿一双舒服的鞋没什么两样
（你也应该穿舒服的鞋）。
对有听觉障碍的人来说，它还能带来很大差别。

\subsection*{镜像学习者的环境}

你可能用花哨的 Unix shell 提示符、
开发环境的自定义配色、
或一大堆键盘快捷键把自己的环境定制过了。
你的学习者这些都没有，
所以要尽量创建一个和他们\emph{实际}拥有的东西一致的环境。
有些老师会在自己的笔记本上另建一个极简账户，
或者如果用的是 Scratch、GitHub 这类在线服务，就单建一个只用于教学的账户。
这么做还有助于防止你昨天为工作安装的软件包
弄坏你今天早上要讲的课。

\subsection*{明智地使用屏幕}

为了让后排的人能看清，
你通常需要把字体放大很多，
这意味着屏幕上能放的东西会比你习惯的少得多。
在很多情况下，你会被限制到 60—70 列、20—30 行，
就像把一台 21 世纪的超级计算机
当成 1980 年代初的哑终端来用。

为了应对这一点，
把你用来讲课的窗口最大化，
然后请所有人用竖大拇指或倒竖大拇指表示它是否读得清。
用浅色背景上的黑色字体，
而不要深色背景上的浅色字体——浅色底比纯白刺眼程度低一些。

也要注意房间里的灯光：
不该全黑，而且你的投影幕布上或正上方不应有直射光。
留几分钟让学习者挪动桌子，
以便能看清楚。

当投影幕布的下沿和学习者的头部一样高时，
后排的人就看不到画面下面那部分。
你可以把窗口的下沿抬高来补偿，
但这样你打字的可用空间就更少了。

如果你能弄到第二台投影仪和幕布，
就用上它：
多出来的空间能让你把代码显示在一侧，
把它的输出或行为显示在另一侧。
如果第二块屏需要单独一台电脑，
就请一位助手来操控它，
而不要在两台键盘之间跳来跳去。

最后，
如果你在控制台窗口里讲 Unix shell 之类的东西，
那就很重要的一点是，要告诉大家你什么时候运行了控制台内的文本编辑器、
什么时候回到了控制台提示符。
大多数新手从没见过一个窗口这样“一人分饰多角”，
很快就会搞不清
你什么时候在和 shell 交互、
什么时候在编辑器里打字、
什么时候在处理 Python 或其他语言的交互式提示符。
你可以用单独一个窗口来编辑，从而避开这个问题；
如果这么做，
一定要在学习者面前说明你正在把焦点从一个窗口切到另一个。

\begin{aside}{无障碍辅助对每个人都有帮助}
  像 \hreffoot{https://boinx.com/mousepose/overview/}{Mouseposé}（Mac 版）
  和 \hreffoot{http://www.pointerfocus.com/}{PointerFocus}（Windows 版）这样的工具
  能在屏幕上高亮鼠标光标的位置，
  而像 \hreffoot{https://www.techsmith.com/video-editor.html}{Camtasia} 这样的录屏工具
  以及 \hreffoot{https://github.com/keycastr/keycastr}{KeyCastr} 这样的独立应用
  会把你敲的 Tab、Control-J 这类看不见的键回显出来。
  刚开始用时，这些工具可能有点烦人，
  但能帮你的学习者弄清你在做什么。
\end{aside}

\subsection*{双设备}

现在有些人在讲课时用两台设备：
一台接入投影仪、给学习者看的笔记本电脑，
外加一台平板，好让自己能看自己的笔记
以及学习者在记的笔记（\secref{s:classroom-notetaking}）。
这比在多个虚拟桌面之间来回切换更可靠，
不过打印出来的讲义仍然是最可靠的备用技术。

\subsection*{早点画，经常画}

示意图总是个好主意。
我有时会准备一整套幻灯片放满提前画好的图，
但一步步把图画出来有助于保持（\secref{s:architecture-brain}），
也让你可以即兴发挥。

\subsection*{避免分心}

关掉笔记本电脑上你可能开启的任何通知，
尤其是来自社交媒体的。
看着屏幕上一闪而过的消息，既会让你分心，也会让学习者分心，
而且当冒出一条你不愿让别人看到的消息时，会让人很尴尬。
再说一遍，
你也许想在电脑上另建一个账户，
完全不用配置电子邮件或其他工具。

\subsection*{熟悉材料之后再即兴发挥}

第一次讲一门课时，
要相当严格地照着你自己拟定或借来的课程计划走。
你可能会很想偏离材料，
因为你想展示一个小窍门，或者演示另一种做法，
但你很有可能会遇到意料之外的事，
花掉的时间会超出你的预算。

不过，一旦你对材料更熟悉了，
你就可以、也应该根据学习者的背景、
他们在课上的问题、
以及你个人觉得最有趣的东西来即兴发挥。
这就像弹一首新歌：
最初几次照着谱子弹，
等你对旋律和和弦变化都熟了，
就可以开始打上自己的印记了。

当你想用一个新东西时，
提前\emph{用你上课要用的那台电脑}把它走一遍：
在高中 WiFi 上当着无聊的十六岁孩子的面安装几百兆的软件，
是你永远不想去经历的。

\begin{aside}{直接教学}
  \gref{g:direct-instruction}{直接教学}是一种教学方法，
  围绕精细设计的课程、通过规定脚本进行讲授。
  它更像演员背台词，而不像我们推荐的即兴发挥式方法。
  \cite{Stoc2018}发现直接教学有统计上显著的正面效果，
  尽管它有时因为机械而被批评。
  我更喜欢即兴发挥，因为直接教学需要的前期准备
  超过大多数散养式学习小组所能负担的。
\end{aside}

\subsection*{偶尔面向屏幕}

在走过一段代码或画图时，
偶尔面向投影屏幕是没问题的：
\emph{不}看一屋子都盯着你的人，
有助于降低你的焦虑水平，也给你一点时间想想接下来该说什么。

但一次不要这样做超过几秒钟。
一条好的经验法则是把投影屏幕当成你的一个学习者：
如果你盯着屏幕的时间长到放在任何人身上都会让人觉得不自在，
那就是时候转回身、重新面对全班了。

\subsection*{缺点}

现场编程确实有一些缺点，
但它们都可以靠一点点练习来避免或绕过。
如果你发现自己犯的打字小错太多，
就每天留出五分钟练打字：
这对你日常工作也有帮助。
如果你觉得花太多时间翻看课程笔记，
就把笔记拆成更小的块，
让自己每次只需想一小步。

如果你觉得自己花了太多时间敲库的 import 语句、
类声明头
和其他样板代码，
就给你自己和学习者一些骨架代码作为起点（\secref{s:classroom-blank}）。
这么做还能降低他们的认知负荷，
因为它会把注意力聚焦到你想让他们关注的地方。

\seclbl{课例研究}{s:performance-jugyokenkyu}

从政客到研究者，再到老师自己，
教育改革者设计过各种制度，
想找出并提拔教得好的人、
淘汰教不好的人。
但「有些人生来就是老师」这个假设是错的：
相反，
和任何其他表演艺术一样，
教得更好的关键是练习和协作。
正如\cite{Gree2014}所解释的，
日本的做法叫做\gref{g:jugyokenkyu}{授業研究}，
意思是「课例研究」：

\begin{quote}

  为了毕业，
  [日本的]教育学专业学生不仅要观摩分配给他的师傅教师上课，
  还要有效地取代他：
  先作为观察者进驻他的课堂，
  到第三周时，
  则成为晃晃悠悠{\ldots}的、近似老师本人的样子。
  这运作起来像一种教学接力。
  每个实习生认领一个科目，
  备好五天的课{\ldots}[然后]各上其中一天。
  要传好这根接力棒，
  你就得在每个科目里各上一天课：
  你备过的那门，以及你没备过的另外四门{\ldots}
  而且你得在师傅教师的眼皮底下做到。
  事后，所有人——师傅教师、大学生们，
  有时还包括另一位外部观察者——会围坐在一张正式的桌子旁
  谈他们看到的东西。

\end{quote}

把作品放到显微镜下改进，这种做法在
\hreffoot{https://en.wikipedia.org/wiki/W.\_Edwards\_Deming}{制造业}和音乐等形形色色的领域里都很常见。
例如，一位职业音乐家在演奏《Body and Soul》或《Smells Like Teen Spirit》之前，
会仔细剖析它的半打录音版本。
他们也期望在练习中和演出后从同行音乐家那里得到反馈。

但在大多数英语国家，持续反馈并不是教学文化的一部分。
在那里，
课堂上发生的事就留在课堂上：
老师不会定期互相听课，
所以也无法借用彼此的好点子。
老师可能从同事、
学校董事会或教科书出版商那里拿到课程计划和作业，
或者在网上上几门 \gref{g:mooc}{MOOC}，
但每个人都得自己琢磨
如何在特定的教室里、为特定的学习者讲好特定的课。
对于参与课后工作坊和集训营的志愿者及其他散养式老师来说，这一点尤其如此。

写出新技巧，
以及上\grefdex{g:demonstration-lesson}{示范课}{示范课}
（由一人给真实学习者授课，其他老师观摩），
都不是解决办法。
例如，
\cite{Finc2007,Finc2012}发现，在分析的 99 个变革故事中，
老师只在三个案例里主动寻找新做法或新材料，
只在八个案例里查阅了已发表的材料。
大多数变革都是本地发生的，
没有外部来源的输入，
或者只涉及与其他教育者的个人互动。
\cite{Bark2015}发现了类似的情况：

\begin{quote}

  采纳不是一种「理性行动」{\ldots}而是
  在社会语境中作出的一系列迭代决策，
  依赖规范传统、社会提示，
  以及情绪性或直觉性的过程{\ldots}
  教师不太可能把教育研究发现
  作为采纳决策的依据{\ldots}
  学生的正面反馈会被教师当作强有力的证据，
  证明他们应该继续某种做法。

\end{quote}

\emph{课例研究}之所以有效，是因为它最大化老师之间非预期知识传递的机会：
某人本打算演示 X，
但看他演示时，
听众实际上也（或反而）学到了 Y。
例如，
一位老师可能想给学习者展示如何在文本文件里搜索电子邮件地址，
但听众带走的可能是几个新的键盘快捷键。

\seclbl{给予和获得关于教学的反馈}{s:performance-feedback}

观察别人对你有帮助，
给别人反馈对他们有帮助，
但接受反馈可能很难，
尤其是负面反馈（\figref{f:performance-feedback-feelings}）。

\figimg{figures/deathbulge-jerk.jpg}{反馈引发的感受（版权所有 © Deathbulge 2013）}{f:performance-feedback-feelings}{}

当双方对什么在范围内、什么不在范围内，
以及评论该怎么措辞有共同预期时，
反馈就更容易给、也更容易收。
如果你是请求反馈的人：\index{反馈!获得}

\begin{description}

\item[主动请求反馈。]
  主动去要反馈，比不情愿地收到反馈要好。

\item[自己选问题，]
  也就是请求具体的反馈。
  让人回答「你觉得怎么样？」
  比回答
  「我讲得太快了吗？」
  或
  「这节课里有哪一件事我应该继续保持？」
  要难得多。
  像这样引导反馈，对你自己也更有帮助。
  一次只设法改一件事，
  总比什么都改、指望会变好要强。
  把反馈引导到你选择要改进的地方，能帮你保持专注，
  这反过来又增加了你看到进步的可能性。

\item[用一位反馈翻译员。]
  让另一个人通读所有反馈，给你一份总结。
  听到「有几个人觉得你可以稍微快一点」，
  通常比读到好几张都写着「这太慢了」
  或「这很无聊」的条子要容易接受。

\item[对自己宽容一点。]
  我们很多人对自己非常苛刻，
  所以在\emph{得到}别人反馈之前，先写下我们自己对自己的看法，
  总是有帮助的。
  这让我们能把自己对表现的看法
  和别人的看法作比较，
  进而更准确地校准前者。
  例如，
  人们常以为自己「嗯」「呃」说得太多，
  而听众其实根本没注意到。
  得到一次这样的反馈，能让老师在下次有同样感觉时调整自我评价。

\end{description}

\noindent
你也可以更有效地给别人反馈：\index{反馈!给予}

\begin{description}

\item[互动。]
  盯着别人看是让他们不自在的好办法，
  所以如果你想就某人平时怎么教给反馈，
  就得先让他们放松。
  像真正的学习者那样和他们互动，是个好办法，
  所以可以提问，或者（假装）跟着他们的例子一起敲。
  如果你是一个小组的一员，
  就安排一两个人扮演学习者，
  其他人做记录。

\item[平衡正面和负面反馈。]
  由一句表扬、一句批评、
  再一句表扬组成的「三明治」，
  很快就会让人厌烦，
  但告诉别人该继续保持什么、以及该改变什么，都很重要\footnote{
    有一阵子，
    我太担心音准，
    结果完全失去了节奏感。
  }。

\item[做记录。]
  如果展示超过几秒钟，
  你就记不住注意到的所有东西，
  更肯定不会记得每件事你注意到了多少次。
  某件事第一次出现时记一笔，
  之后它再出现就加一个勾，
  这样你就能按频率来整理你的反馈。

\end{description}

当你手头有某种评分标准时，做记录会更高效，
免得你一边慌着写观察，一边被观察的人还在说话。
对来自小组的自由评论来说，
最简单的评分标准是一张 2x2 的网格，
纵轴标「哪些做得好」和「哪些可以改进」，
横轴标「内容」（说了什么）
和「呈现」（怎么说的）。
观察者在看示范时把评论写在便利贴上，
然后把它们贴到白板上画好的网格对应的四个象限里
（\figref{f:performance-rubric}）。

\figpdf{figures/2x2-rubric.pdf}{教学评分标准}{f:performance-rubric}{}

\begin{aside}{评分标准与问题配额}
  \secref{s:checklists-teacheval}里有一份评分标准样例，
  用来评估 5—10 分钟的编程教学。
  它大体按问题可能出现的顺序排列，
  例如关于引入的问题排在关于收尾的问题之前。

  这类评分标准往往会随时间不断膨胀，因为人们总想加东西。
  让它们保持可控的一个好办法，是坚持总长度不变：
  如果有人想加一个问题，
  就得找出一条不那么重要、可以删掉的问题。
\end{aside}

如果你对给予和获得反馈感兴趣，
\cite{Gorm2014}有一些好建议，
可以用来把同伴互评变成你教学的常规部分；
而~\cite{Gawa2011}探讨了拥有一位教练的价值。

\begin{aside}{工作室课程}
  建筑学院常设有工作室课程，\index{工作室课程}
  学生在其中解决小的设计问题，
  并当场从同伴那里得到反馈。
  当老师对同伴的评论再作评论时，这类课最有效，
  这样参与者不仅学会怎么造房子，
  还学会怎么给和收反馈~\cite{Scho1984}。
  音乐里的 master class 也有类似作用，
  我发现给反馈再给反馈，
  也能帮人们改进教学。
\end{aside}

\seclbl{如何练习呈现}{s:performance-practice}

改进当面授课最好的办法，
是看你自己的表现：

\begin{itemize}

\item
  三人一组活动。

\item
  每个人轮流扮演老师、听众和摄像师。
  老师有 2 分钟解释某样东西。
  扮听众的人负责专注倾听，
  而摄像师用手机或其他手持设备把这段录下来。

\item
  所有人都讲完之后，
  全组一起看这些视频。
  每个人对三段视频都给出反馈，
  也就是人们既评论自己，也评论别人。

\item
  视频讨论完后就删掉。
  （很多人对自出现在网上感到不自在，这是有道理的。）

\item
  最后，
  全班重新集合，
  把所有反馈汇总到上面描述过的那种共享 2x2 网格里，
  \emph{不要}说明每条反馈是关于谁的。

\end{itemize}

要让这个练习效果好：

\begin{itemize}

\item
  把三段视频都录完，然后一起看。
  如果循环是「教—评—教—评」，
  最后一个讲的人总是不够时间
  （有时是故意的）。
  把所有评议都放在所有授课之后做，
  也有助于在两者之间拉开一点距离，
  让这个练习稍微没那么难熬。

\item
  上课一开始就告诉大家他们会被要求讲点东西，
  好让他们有时间选题。
  过早告诉他们可能适得其反，
  因为有些人会为「该准备到什么程度」而发愁。

\item
  各组必须在空间上分开，以减少录音之间的声音串扰。
  实际上，
  这意味着在一间普通教室里放 2—3 组，
  其余的人用附近的休息区、咖啡室、办公室，
  或者（有一次）用清洁工的储物间。

\item
  人们必须既给自己、也给别人反馈，
  这样才能把自己对自己教学的印象
  同别人的印象校准起来。
  大多数人对自己比应该的更苛刻，
  意识到这一点对他们很重要。

\end{itemize}

宣布这个练习时，往往招来一片哀嚎和忐忑，
因为很少有人喜欢看到或听到自己。
然而，
同样是这些人，却一直把它评为教师工作坊里最有价值的部分之一。
它也是很好的协同教学准备（\secref{s:classroom-together}）：\index{协同教学}
老师们发现，如果已经练习过互相给非正式反馈、
并有一套共同的评分标准来设定预期，
再互相给反馈就容易得多。

说到评分标准：
等全班把所有反馈都放到共享网格上之后，
挑出少量正面和负面评论，
把它们写成一份检查清单，
然后让大家再做一遍这个练习。
大多数人第二遍会自在得多，
而按他们自己认定重要的东西来评价，
会增强他们的自主感（\chapref{s:motivation}）。

\begin{aside}{小动作}
  我们都有紧张的习惯：
  上台时会比平常说得更快、音调更高，
  会玩头发，
  或者掰手指关节。
  赌徒把这些叫做「破绽」（tells），
  而人们往往没有意识到，当自己其实不知道某个问题的答案时，
  会来回踱步、
  盯着自己的鞋、
  或者把口袋里的零钱弄得叮当响。

  你没法完全去掉这些小动作，
  一味想改掉反而可能让你为此着迷。
  更好的策略是设法把它们替换掉——例如，
  训练自己在紧张时用脚趾在鞋里蜷一蜷，
  而不是用小指去掏耳朵。
\end{aside}

\seclbl{练习}{s:performance-exercises}

\exercise{对糟糕的教学给反馈}{全班}{20}

作为一个集体，
看\hreffoot{https://www.youtube.com/watch?v=jxgMVwQamO0}{这段糟糕教学的视频}，
从两个维度给反馈：
正面 vs.\ 负面，内容 vs.\ 呈现。
让班上每个人往白板上或共享笔记里的 2x2 网格里加一条，
不要和别人重复。
别人看到了哪些你漏掉的？
他们有什么看法是你强烈同意或不同意的？

\exercise{练习给反馈}{小组}{45}

用\secref{s:performance-practice}里描述的过程，
以三人一组练习讲课，
并汇总反馈。

\exercise{差与好}{全班}{20}

看\hreffoot{https://youtu.be/bXxBeNkKmJE}{做得差的现场编程}
和\hreffoot{https://youtu.be/SkPmwe\_WjeY}{做得好的现场编程}
这两段视频，用惯常的 2x2 网格总结你对两者差异的反馈。
第二轮教学为什么比第一轮好？
第一轮里有没有什么比第二轮做得更好的地方？

\exercise{先看，再做}{两人一组}{30}

用现场编程给一位同学讲 3—4 分钟的课，
然后交换，看对方为你现场编程。
不必费劲去录这些片段——用手持设备很难同时拍到人和屏幕——
但用你之前一样的办法给反馈。
提前向你的同侪说明你要讲什么，
以及你讲给的学习者预期已经熟悉哪些内容。

\begin{itemize}

\item
  相比站起来讲课，现场编程感觉有什么不同？
  哪些更容易或更难？

\item
  你犯了错吗？
  如果有，你是怎么处理的？

\item
  你是一边说话一边打字，还是交替进行？

\item
  你多久指一次屏幕？
  多久用鼠标高亮一次？

\item
  下次你打算继续保持哪些做法？
  打算换哪些做法？

\end{itemize}

\exercise{小动作}{小组}{15}

\begin{enumerate}

\item
  记下你认为自己有哪些小动作，
  但不要和其他人分享。

\item
  讲一节短课（2—3 分钟）。

\item
  问你的听众，他们觉得你怎么暴露了紧张。
  他们的清单和你的一样吗？

\end{enumerate}

\exercise{教学提示}{小组}{15}

\hreffoot{http://csteachingtips.org/}{CS Teaching Tips} 网站
有大量关于计算教学的实用提示，
还有一套可下载的提示单。
以小组为单位过一遍他们的提示单，把每条提示分类：
你是总用、偶尔用、还是从不用。
你的实践和同伴的实践在哪里不同？
有没有哪些提示你强烈不同意、或认为不会有效？

\section*{回顾}

\figpdfhere{figures/conceptmap-feedback.pdf}{概念：反馈}{f:performance-feedback}{}
